% ============================================================
%  RAPPORT DE STAGE DE FIN D'ÉTUDES
%  Plateforme web de gestion des formations — STEG
%
%  Compiler avec pdfLaTeX (recommandé : Overleaf)
%  Les éléments à personnaliser sont marqués [[ ... ]]
% ============================================================
\documentclass[12pt,a4paper]{report}

% ---------- Encodage & langue ----------
\usepackage[utf8]{inputenc}
\usepackage[T1]{fontenc}
\usepackage{lmodern}
\usepackage[french]{babel}

% ---------- Mise en page ----------
\usepackage[a4paper,left=2.5cm,right=2.5cm,top=3cm,bottom=3cm]{geometry}
\usepackage{setspace}
\onehalfspacing
\usepackage{fancyhdr}
\pagestyle{fancy}
\fancyhf{}
\fancyhead[L]{\small\nouppercase{\leftmark}}
\fancyfoot[C]{\thepage}
\renewcommand{\headrulewidth}{0.4pt}
\usepackage{microtype}

% ---------- Tableaux & figures ----------
\usepackage{graphicx}
\graphicspath{{figures/}}
\usepackage{booktabs}
\usepackage{tabularx}
\usepackage{longtable}
\usepackage{multirow}
\usepackage{array}
\usepackage{float}
\usepackage{caption}
\captionsetup{font=small,labelfont=bf}
\renewcommand{\arraystretch}{1.35}
\newcolumntype{C}[1]{>{\centering\arraybackslash}p{#1}}

% ---------- Divers ----------
\usepackage{enumitem}
\setlist{itemsep=2pt,topsep=4pt}
\usepackage{xcolor}
\usepackage{tcolorbox}
\usepackage{pgfgantt}
\usepackage{listings}
\lstset{basicstyle=\ttfamily\small,frame=single,keywordstyle=\color{blue},
  commentstyle=\color{gray},breaklines=true,showstringspaces=false}
\usepackage{pdfpages}
\usepackage[colorlinks=true,linkcolor=black,citecolor=black,urlcolor=blue]{hyperref}
\hypersetup{
  pdftitle={Rapport de stage — Plateforme de gestion des formations STEG},
  pdfauthor={[[Prénom NOM]]}
}

% ============================================================
\begin{document}

% ------------------------------------------------------------
%  PAGE DE GARDE
% ------------------------------------------------------------
\begin{titlepage}
\begin{center}

{\large [[Université / École — ex. Université de Tunis El Manar]]\\
[[École Nationale d'Ingénieurs / Faculté / Institut Supérieur]]\\[0.3cm]}
\rule{\textwidth}{1pt}\\[0.5cm]

\textbf{\LARGE Rapport de stage de fin d'études}\\[0.2cm]
Pour l'obtention du diplôme de : \textbf{[[Diplôme — ex. Licence / Ingénieur en Informatique]]}\\[0.8cm]

\includegraphics[width=3cm]{logo-universite.png}\hfill
\includegraphics[width=3cm]{logo-steg.png}\\[0.3cm]
{\footnotesize (logo université) \hfill (logo STEG)}\\[1cm]

\begin{tcolorbox}[colback=blue!5,colframe=blue!40!black,arc=3mm]
\centering\textbf{\Large Conception et réalisation d'une plateforme web\\
de gestion des formations}\\[0.2cm]
\textbf{Étude de cas : Société Tunisienne de l'Électricité et du Gaz (STEG)}
\end{tcolorbox}

\vspace{1cm}
\begin{tabular}{p{7cm} p{7cm}}
\textbf{Élaboré par :} & \textbf{Encadré par :} \\
Lotfi [[NOM]] & [[M. Prénom NOM]] — Encadrant académique \\
\small ([[Classe / Groupe]]) & [[M. Prénom NOM]] — Encadrant professionnel STEG \\
\end{tabular}

\vspace{1.5cm}
\textbf{Année universitaire : [[2025 -- 2026]]}
\end{center}
\end{titlepage}

% ------------------------------------------------------------
%  PAGES PRÉLIMINAIRES
% ------------------------------------------------------------
\pagenumbering{roman}

% --- Dédicace ---
\chapter*{Dédicace}
\addcontentsline{toc}{chapter}{Dédicace}
\begin{flushright}
\itshape [[Texte de dédicace : à mes parents, à ma famille, à mes amis\dots]]
\end{flushright}

% --- Remerciements ---
\chapter*{Remerciements}
\addcontentsline{toc}{chapter}{Remerciements}
[[Remercier : l'encadrant professionnel STEG, l'encadrant académique,
l'équipe du service formation de la STEG, et toute personne ayant contribué
au bon déroulement du stage.]]

% --- Résumés ---
\chapter*{Résumé}
\addcontentsline{toc}{chapter}{Résumé}
Ce rapport présente la conception et la réalisation d'une plateforme web
de gestion des formations destinée à la Société Tunisienne de l'Électricité
et du Gaz (STEG). Le système couvre l'ensemble du cycle de formation :
gestion du catalogue, planification des sessions, gestion des participants
et des employés, suivi des présences, évaluation, génération des documents
officiels (conventions, contrats de formateurs, certificats), signature
électronique, ainsi qu'un assistant intelligent basé sur la technique du
\emph{Retrieval-Augmented Generation} (RAG) pour l'exploitation automatique
des documents de formation.

Le projet a été mené selon la méthodologie \textbf{Scrum} et a abouti à deux
releases fonctionnelles. La solution a été développée avec
\textbf{NestJS} et \textbf{TypeORM} côté serveur, \textbf{React} avec
\textbf{TypeScript} côté client, et une base de données \textbf{MySQL}.

\textbf{Mots-clés :} gestion des formations, STEG, NestJS, React,
TypeScript, Scrum, RAG, signature électronique.

\chapter*{Abstract}
\addcontentsline{toc}{chapter}{Abstract}
This report presents the design and implementation of a web-based training
management platform for the Tunisian Company of Electricity and Gas (STEG).
The system covers the whole training life cycle: catalog management, session
planning, participant management, attendance tracking, evaluations, official
document generation (conventions, trainer contracts, certificates),
e-signatures, and an AI assistant based on Retrieval-Augmented Generation
(RAG). The project followed the \textbf{Scrum} methodology and was delivered
in two releases, using \textbf{NestJS}, \textbf{React}, \textbf{TypeScript}
and \textbf{MySQL}.

\textbf{Keywords:} training management, STEG, NestJS, React, TypeScript,
Scrum, RAG, e-signature.

% NB : le résumé en arabe doit être composé avec XeLaTeX/LuaLaTeX
% (babel arabic ou polyglossia), ou inséré comme image :
\chapter*{ملخص}
\addcontentsline{toc}{chapter}{Résumé (arabe)}
% \begin{otherlanguage}{arabic} ... \end{otherlanguage}  % (nécessite XeLaTeX)
\begin{center}
\fbox{\parbox[c][5cm][c]{0.8\textwidth}{\centering
Insérer ici le résumé en arabe\\ (image ou compilation XeLaTeX)}}
\end{center}

% --- Table des matières, figures, tableaux ---
\tableofcontents
\listoffigures
\listoftables

% --- Glossaire ---
\chapter*{Glossaire}
\addcontentsline{toc}{chapter}{Glossaire}
\begin{description}[style=nextline]
  \item[API] \emph{Application Programming Interface} : interface de
    programmation permettant la communication entre applications.
  \item[CNFCPP] Fonds de compensation pour la formation professionnelle continue.
  \item[CRUD] \emph{Create, Read, Update, Delete} : opérations de base sur les données.
  \item[JWT] \emph{JSON Web Token} : standard de jetons d'authentification.
  \item[KPI] \emph{Key Performance Indicator} : indicateur clé de performance.
  \item[RAG] \emph{Retrieval-Augmented Generation} : génération augmentée par
    recherche dans une base documentaire.
  \item[REST] \emph{REpresentational State Transfer} : style d'architecture d'API web.
  \item[STEG] Société Tunisienne de l'Électricité et du Gaz.
  \item[UML] \emph{Unified Modeling Language} : langage de modélisation orienté objet.
  \item[User Story] Expression d'un besoin fonctionnel du point de vue de l'utilisateur.
\end{description}

% ============================================================
%  INTRODUCTION GÉNÉRALE
% ============================================================
\chapter*{Introduction générale}
\addcontentsline{toc}{chapter}{Introduction générale}

Dans un environnement concurrentiel en constante évolution, la compétence
du personnel constitue le premier facteur de performance d'une entreprise.
La Société Tunisienne de l'Électricité et du Gaz (STEG) accorde une importance
particulière à la formation continue de ses employés afin d'accompagner
l'évolution technologique de ses métiers. Toutefois, la gestion de ces
formations reste largement manuelle : planification des sessions, suivi des
présences, production des documents officiels et archivage sont réalisés
avec des outils dispersés (fichiers Word, tableurs, e-mails), ce qui engendre
perte de temps, erreurs et difficultés de traçabilité.

Ce rapport présente la conception et la réalisation d'une \textbf{plateforme
web de gestion des formations} destinée à la STEG. La solution couvre le cycle
complet de la formation : catalogue des formations, planification des sessions,
inscriptions, gestion des présences, évaluations, génération et signature
électronique des documents officiels, tableaux de bord de pilotage, et un
assistant intelligent d'exploitation des documents de formation.

Le présent rapport est organisé en quatre chapitres :
\begin{itemize}
  \item \textbf{Chapitre 1} : le cadre du projet — présentation de
    l'organisme d'accueil, problématique, solution proposée et méthodologie de travail.
  \item \textbf{Chapitre 2} : analyse et spécification des besoins —
    spécification fonctionnelle (backlogs) et non fonctionnelle.
  \item \textbf{Chapitre 3} : conception — diagrammes de classes, diagrammes
    de séquence et modèle de données.
  \item \textbf{Chapitre 4} : réalisation — environnement technique,
    architecture, interfaces réalisées et tests.
\end{itemize}
Une conclusion générale clôture le rapport et ouvre des perspectives
d'évolution.

% ============================================================
%  CHAPITRE 1 — LE CADRE DU PROJET
% ============================================================
\chapter{Le cadre du projet}
\pagenumbering{arabic}\setcounter{page}{1}

\section{Présentation de l'organisme d'accueil}
\subsection{Fiche signalétique}
\begin{table}[H]
\centering
\begin{tabular}{p{4.5cm} p{8.5cm}}
\toprule
\textbf{Raison sociale} & Société Tunisienne de l'Électricité et du Gaz (STEG) \\
\textbf{Secteur d'activité} & Production, transport et distribution d'électricité et de gaz \\
\textbf{Création} & 1962 \\
\textbf{Statut} & Entreprise publique tunisienne \\
\textbf{Siège} & Tunis \\
\textbf{Effectif} & [[~16 000 agents]] \\
\bottomrule
\end{tabular}
\caption{Fiche signalétique de la STEG}
\end{table}

\subsection{Organisation et service d'accueil}
[[Décrire brièvement la structure de la STEG et présenter le service d'accueil
(direction de la formation / centre de formation), ses missions et son
fonctionnement actuel.]]

\section{Contexte et problématique}
\subsection{Étude de l'existant}
[[Décrire le processus actuel de gestion des formations au sein de la STEG :
dossiers papier, tableurs Excel, échanges par e-mail, absence d'historique
centralisé des présences et des évaluations.]]

\subsection{Critique de l'existant}
\begin{itemize}
  \item Dispersion des données entre plusieurs fichiers et intervenants ;
  \item Production manuelle des documents officiels (conventions,
    contrats, certificats), longue et source d'erreurs ;
  \item Absence de traçabilité des présences et des évaluations ;
  \item Difficulté à produire des statistiques fiables pour le pilotage ;
  \item Documents de formation (supports PDF) difficilement exploitables.
\end{itemize}

\begin{tcolorbox}[colback=blue!5,colframe=blue!40!black,title=Problématique]
\itshape Comment concevoir une plateforme web centralisée, sécurisée et
intuitive, permettant de gérer l'intégralité du cycle de formation au sein de
la STEG — de la planification à la délivrance des documents officiels — tout
en automatisant la production documentaire et en offrant une visibilité
temps réel sur les indicateurs de formation ?
\end{tcolorbox}

\section{Solution proposée}
Développement d'une application web composée de :
\begin{itemize}
  \item une \textbf{application front-end} (React) avec des espaces dédiés
    selon les rôles : administrateur, formateur, participant ;
  \item une \textbf{API REST sécurisée} (NestJS) gérant les formations,
    sessions, inscriptions, présences, évaluations et documents ;
  \item une \textbf{signature électronique} des documents officiels ;
  \item un \textbf{assistant intelligent} (résumés automatiques et
    question/réponse) s'appuyant sur les supports de formation (\emph{RAG}).
\end{itemize}

\section{Méthodologie de travail}
\subsection{Méthodologie Scrum}
Le projet a été conduit selon la méthodologie agile \textbf{Scrum} :
\begin{itemize}
  \item \textbf{Product Owner} : encadrant professionnel STEG ;
  \item \textbf{Scrum Master / équipe de développement} : stagiaire ;
  \item \textbf{Backlog du produit} : ensemble des fonctionnalités priorisées
    (cf.~chapitre~2) ;
  \item \textbf{Sprints} : cycles de deux semaines avec revue et rétrospective ;
  \item \textbf{Deux releases} : la première couvre la gestion de base des
    formations, la seconde ajoute les documents, la traçabilité et l'IA.
\end{itemize}

\subsection{Planification du projet}
\begin{figure}[H]
\centering
\begin{ganttchart}[
    hgrid,vgrid,
    x unit=0.5cm,
    y unit title=0.75cm,
    y unit chart=0.6cm,
    bar/.append style={fill=blue!60},
    bar height=0.6,
    group/.append style={fill=black!70},
    group height=0.4,
    milestone/.append style={fill=red}
]{1}{16}
\gantttitle{Planification du stage (4 mois)}{16}\\
\gantttitle{M1}{4}\gantttitle{M2}{4}\gantttitle{M3}{4}\gantttitle{M4}{4}\\
\ganttgroup{Étude de l'existant \& cadre du projet}{1}{2}\\
\ganttbar{Spécification des besoins}{2}{4}\\
\ganttbar{Conception (UML, BD)}{4}{6}\\
\ganttgroup{Release 1 : gestion des formations}{6}{10}\\
\ganttbar{Sprint 1 : authentification \& comptes}{6}{7}\\
\ganttbar{Sprint 2 : formations, sessions, présences}{8}{10}\\
\ganttgroup{Release 2 : documents, KPI \& IA}{11}{15}\\
\ganttbar{Sprint 3 : documents \& signatures}{11}{13}\\
\ganttbar{Sprint 4 : KPI \& assistant IA}{13}{15}\\
\ganttbar{Rédaction du rapport}{14}{16}\\
\ganttmilestone{Soutenance}{16}
\end{ganttchart}
\caption{Diagramme de Gantt du projet}
\end{figure}

\section{Environnement de travail}
\begin{table}[H]
\centering
\begin{tabular}{p{4.5cm} p{8.5cm}}
\toprule
\textbf{Matériel} & PC portable [[i5/i7, 16 Go RAM]] \\
\textbf{Système d'exploitation} & Windows 11 \\
\textbf{IDE} & Visual Studio Code \\
\textbf{Gestion de versions} & Git / GitHub \\
\textbf{Modélisation} & draw.io (diagrammes UML) \\
\textbf{Tests d'API} & Postman \\
\textbf{Base de données} & MySQL + phpMyAdmin \\
\bottomrule
\end{tabular}
\caption{Environnement de travail}
\end{table}

% ============================================================
%  CHAPITRE 2 — ANALYSE ET SPÉCIFICATION DES BESOINS
% ============================================================
\chapter{Analyse et spécification des besoins}

\section{Identification des acteurs}
\begin{table}[H]
\centering
\begin{tabular}{p{3.5cm} p{9.5cm}}
\toprule
\textbf{Acteur} & \textbf{Description} \\
\midrule
Administrateur & Gère les comptes, les formations, les sessions, les
documents officiels et supervise les indicateurs de formation. \\
Formateur & Anime les formations, consulte son planning, gère les présences
et consulte les évaluations de ses sessions. \\
Participant (employé) & Consulte le catalogue, s'inscrit aux formations,
assiste aux sessions, évalue les formations et télécharge ses certificats. \\
Entreprise (cabinet) & [[Si utilisé : structure rattachée regroupant des
employés, avec un compte de gestion dédié.]] \\
Assistant IA & Acteur système fournissant résumés et réponses sur les
documents de formation. \\
\bottomrule
\end{tabular}
\caption{Identification des acteurs}
\end{table}

\section{Spécification fonctionnelle}
\subsection{Backlog du produit}
\begin{table}[H]
\small
\begin{tabularx}{\textwidth}{C{1cm} C{2.2cm} X C{1.8cm}}
\toprule
\textbf{ID} & \textbf{En tant que\dots} & \textbf{Je veux\dots} & \textbf{Priorité} \\
\midrule
US1 & administrateur & m'authentifier et gérer les comptes (administrateurs, formateurs, participants) & Haute \\
US2 & administrateur & gérer les formations : création, modification, publication au catalogue & Haute \\
US3 & administrateur & planifier des sessions et affecter formateurs et salles & Haute \\
US4 & participant & consulter le catalogue et m'inscrire à une formation & Haute \\
US5 & formateur & gérer les présences des participants aux sessions & Moyenne \\
US6 & participant & évaluer les formations et consulter mes résultats & Moyenne \\
US7 & administrateur & générer les documents officiels (conventions, contrats de formateurs, certificats) en PDF & Haute \\
US8 & participant & signer électroniquement les documents & Haute \\
US9 & administrateur & archiver et télécharger les documents signés & Moyenne \\
US10 & administrateur & consulter un tableau de bord d'indicateurs (taux de participation, moyennes d'évaluation) & Moyenne \\
US11 & utilisateur & obtenir des résumés automatiques et poser des questions sur les documents de formation (IA) & Basse \\
\bottomrule
\end{tabularx}
\caption{Backlog du produit}
\end{table}

\subsection{Backlog du Release 1 — gestion des formations}
\begin{table}[H]
\small
\begin{tabularx}{\textwidth}{C{1cm} C{2.2cm} X C{1.6cm} C{1.4cm}}
\toprule
\textbf{ID} & \textbf{En tant que\dots} & \textbf{Je veux\dots} & \textbf{Priorité} & \textbf{Sprint} \\
\midrule
US1 & administrateur & m'authentifier et gérer les comptes & Haute & S1 \\
US2 & administrateur & gérer les formations et le catalogue & Haute & S2 \\
US3 & administrateur & planifier les sessions & Haute & S2 \\
US4 & participant & m'inscrire à une formation & Haute & S2 \\
US5 & formateur & gérer les présences & Moyenne & S2 \\
US6 & participant & évaluer les formations & Moyenne & S2 \\
\bottomrule
\end{tabularx}
\caption{Backlog du Release 1}
\end{table}

\subsection{Backlog du Release 2 — documents, traçabilité et IA}
\begin{table}[H]
\small
\begin{tabularx}{\textwidth}{C{1cm} C{2.2cm} X C{1.6cm} C{1.4cm}}
\toprule
\textbf{ID} & \textbf{En tant que\dots} & \textbf{Je veux\dots} & \textbf{Priorité} & \textbf{Sprint} \\
\midrule
US7 & administrateur & générer conventions, contrats et certificats en PDF & Haute & S3 \\
US8 & participant & signer électroniquement les documents & Haute & S3 \\
US9 & administrateur & archiver et télécharger les documents signés & Moyenne & S3 \\
US10 & administrateur & consulter le tableau de bord des KPI & Moyenne & S4 \\
US11 & utilisateur & utiliser l'assistant IA sur les documents & Basse & S4 \\
\bottomrule
\end{tabularx}
\caption{Backlog du Release 2}
\end{table}

\subsection{Diagrammes de cas d'utilisation}
% Pour chaque figure : exportez le diagramme draw.io en PNG dans le dossier
% figures/ puis remplacez l'encadré par \includegraphics{...}
\begin{figure}[H]
\centering
\fbox{\parbox[c][8cm][c]{0.85\textwidth}{\centering
\textbf{Insérer ici :} diagramme de cas d'utilisation global\\
(figures/cas-utilisation-global.png)}}
\caption{Diagramme de cas d'utilisation global}
\end{figure}

\begin{figure}[H]
\centering
\fbox{\parbox[c][7cm][c]{0.85\textwidth}{\centering
\textbf{Insérer ici :} cas d'utilisation de l'administrateur\\
(figures/cas-utilisation-admin.png)}}
\caption{Cas d'utilisation de l'administrateur}
\end{figure}

\begin{figure}[H]
\centering
\fbox{\parbox[c][7cm][c]{0.85\textwidth}{\centering
\textbf{Insérer ici :} cas d'utilisation du formateur\\
(figures/cas-utilisation-formateur.png)}}
\caption{Cas d'utilisation du formateur}
\end{figure}

\begin{figure}[H]
\centering
\fbox{\parbox[c][7cm][c]{0.85\textwidth}{\centering
\textbf{Insérer ici :} cas d'utilisation du participant\\
(figures/cas-utilisation-participant.png)}}
\caption{Cas d'utilisation du participant}
\end{figure}

\subsection{Description textuelle des cas d'utilisation}
\begin{table}[H]
\small
\begin{tabularx}{\textwidth}{p{3.5cm} X}
\toprule
\textbf{Cas} & S'authentifier \\
\midrule
Acteur principal & Tout utilisateur \\
Précondition & Posséder un compte actif \\
Scénario nominal &
\begin{enumerate}
  \item L'utilisateur saisit son e-mail et son mot de passe ;
  \item Le système vérifie les identifiants ;
  \item Le système délivre un jeton JWT et redirige selon le rôle.
\end{enumerate} \\
Scénario alternatif & Identifiants invalides : message d'erreur. \\
Postcondition & Session authentifiée. \\
\bottomrule
\end{tabularx}
\caption{Description du cas d'utilisation « S'authentifier »}
\end{table}

\section{Spécification non fonctionnelle}
\begin{itemize}
  \item \textbf{Sécurité} : authentification par JWT, mots de passe hachés
    (bcrypt), contrôle d'accès par rôles (guards) sur chaque route de l'API ;
  \item \textbf{Performance} : temps de réponse des pages inférieur à 2
    secondes, pagination des listes volumineuses ;
  \item \textbf{Ergonomie} : interface responsive, claire et adaptée aux
    profils non techniques ;
  \item \textbf{Traçabilité} : historique des présences, évaluations et
    documents archivés ;
  \item \textbf{Maintenabilité} : architecture modulaire (modules NestJS),
    code typé TypeScript, séparation front-end / back-end ;
  \item \textbf{Extensibilité} : ajout aisé de nouveaux modules (ex. paie
    des vacataires, gestion des salles).
\end{itemize}

% ============================================================
%  CHAPITRE 3 — CONCEPTION
% ============================================================
\chapter{Conception}

\section{Diagrammes de classes}
\begin{figure}[H]
\centering
\fbox{\parbox[c][10cm][c]{0.9\textwidth}{\centering
\textbf{Insérer ici :} diagramme de classes global\\
(figures/classes-global.png)\\[0.3cm]
Principales entités : User, Formation, Session, Inscription, Presence,
Evaluation, SessionDocument, DocumentSigne, Signature, Certificat}}
\caption{Diagramme de classes global}
\end{figure}

\begin{figure}[H]
\centering
\fbox{\parbox[c][8cm][c]{0.9\textwidth}{\centering
\textbf{Insérer ici :} diagramme de classes détaillé du module
Sessions/Présences (figures/classes-sessions.png)}}
\caption{Diagramme de classes du module Sessions}
\end{figure}

\section{Diagrammes de séquence}
\begin{figure}[H]
\centering
\fbox{\parbox[c][8cm][c]{0.9\textwidth}{\centering
\textbf{Insérer ici :} séquence « Inscription à une formation »\\
(figures/sequence-inscription.png)}}
\caption{Diagramme de séquence : inscription à une formation}
\end{figure}

\begin{figure}[H]
\centering
\fbox{\parbox[c][8cm][c]{0.9\textwidth}{\centering
\textbf{Insérer ici :} séquence « Génération et signature du certificat »\\
(figures/sequence-certificat.png)}}
\caption{Diagramme de séquence : génération et signature du certificat}
\end{figure}

\begin{figure}[H]
\centering
\fbox{\parbox[c][8cm][c]{0.9\textwidth}{\centering
\textbf{Insérer ici :} séquence « Question posée à l'assistant IA (RAG) »\\
(figures/sequence-rag.png)}}
\caption{Diagramme de séquence : assistant IA (RAG)}
\end{figure}

\section{Modèle de données}
\subsection{Architecture de la base de données}
La base de données \textbf{MySQL} comprend principalement les tables :
\texttt{users}, \texttt{formations}, \texttt{sessions},
\texttt{inscriptions}, \texttt{presences}, \texttt{evaluations},
\texttt{session\_documents}, \texttt{documents\_signes},
\texttt{signatures} et \texttt{certificats}.

\begin{figure}[H]
\centering
\fbox{\parbox[c][9cm][c]{0.9\textwidth}{\centering
\textbf{Insérer ici :} schéma relationnel de la base de données\\
(figures/modele-donnees.png)}}
\caption{Modèle relationnel de la base de données}
\end{figure}

\subsection{Principales règles de gestion}
\begin{itemize}
  \item Un participant ne peut s'inscrire qu'une seule fois à une session ;
  \item Un certificat est unique par couple (participant, session) ;
  \item Un document signé ne peut plus être modifié après signature ;
  \item Seul un administrateur peut valider une formation pour publication
    au catalogue ;
  \item Les documents officiels générés sont archivés automatiquement.
\end{itemize}

% ============================================================
%  CHAPITRE 4 — RÉALISATION
% ============================================================
\chapter{Réalisation}

\section{Environnement technique}
\begin{table}[H]
\centering
\begin{tabular}{p{3cm} p{10cm}}
\toprule
\textbf{Domaine} & \textbf{Technologies} \\
\midrule
Front-end & React 18, TypeScript, TanStack Router, TanStack Query,
TailwindCSS, shadcn/ui \\
Back-end & NestJS (Node.js), TypeScript \\
ORM & TypeORM \\
Base de données & MySQL \\
Authentification & JWT, bcrypt, guards par rôles \\
Génération PDF & pdfkit \\
IA & Embeddings vectoriels + RAG sur les supports de formation \\
Outils & VS Code, Git/GitHub, Postman, phpMyAdmin \\
\bottomrule
\end{tabular}
\caption{Environnement technique du projet}
\end{table}

\section{Architecture de l'application}
\subsection{Architecture générale}
L'application suit une architecture \textbf{client-serveur} à trois tiers :
\begin{itemize}
  \item \textbf{Présentation} : application React communiquant avec l'API
    REST en JSON ;
  \item \textbf{Métier} : API NestJS organisée en modules
    (auth, users, formations, sessions, inscriptions, presences,
    evaluations, signature, certificate, ai, chatbot) ;
  \item \textbf{Données} : base MySQL manipulée via TypeORM.
\end{itemize}

\begin{figure}[H]
\centering
\fbox{\parbox[c][7cm][c]{0.9\textwidth}{\centering
\textbf{Insérer ici :} schéma de l'architecture trois tiers\\
(figures/architecture.png)}}
\caption{Architecture générale de l'application}
\end{figure}

\subsection{Organisation du back-end}
\begin{lstlisting}[caption={Extrait : organisation modulaire du back-end (NestJS)}]
src/
  common/         # enums, gardes, utilitaires
  config/         # configuration TypeORM
  entities/       # entites TypeORM (User, Formation, Session...)
  modules/
    auth/         # JWT, guards, strategie
    formation/    # CRUD formations + catalogue
    session/      # sessions, participants, presences
    signature/    # conventions, contrats, signatures
    certificate/  # certificats de formation
    ai/           # embeddings + RAG
    chatbot/      # assistant conversationnel
\end{lstlisting}

\section{Interfaces réalisées}
\subsection{Release 1 : gestion des formations}
\begin{figure}[H]
\centering
\fbox{\parbox[c][6cm][c]{0.85\textwidth}{\centering
\textbf{Capture :} page d'authentification (figures/connexion.png)}}
\caption{Interface d'authentification}
\end{figure}

\begin{figure}[H]
\centering
\fbox{\parbox[c][6cm][c]{0.85\textwidth}{\centering
\textbf{Capture :} catalogue des formations (figures/catalogue.png)}}
\caption{Catalogue des formations}
\end{figure}

\begin{figure}[H]
\centering
\fbox{\parbox[c][6cm][c]{0.85\textwidth}{\centering
\textbf{Capture :} gestion des sessions par l'administrateur
(figures/admin-sessions.png)}}
\caption{Gestion des sessions (back-office)}
\end{figure}

\subsection{Release 2 : documents, KPI et assistant IA}
\begin{figure}[H]
\centering
\fbox{\parbox[c][6cm][c]{0.85\textwidth}{\centering
\textbf{Capture :} signature électronique d'un document
(figures/signature.png)}}
\caption{Signature électronique}
\end{figure}

\begin{figure}[H]
\centering
\fbox{\parbox[c][6cm][c]{0.85\textwidth}{\centering
\textbf{Capture :} tableau de bord des indicateurs KPI
(figures/admin-kpi.png)}}
\caption{Tableau de bord des indicateurs}
\end{figure}

\begin{figure}[H]
\centering
\fbox{\parbox[c][6cm][c]{0.85\textwidth}{\centering
\textbf{Capture :} assistant IA (chatbot) sur les documents de formation
(figures/chatbot.png)}}
\caption{Assistant IA (RAG)}
\end{figure}

\section{Tests et validation}
\subsection{Tests fonctionnels}
Les fonctionnalités ont été validées manuellement et via \textbf{Postman}
pour l'API. Extrait de plan de tests :
\begin{table}[H]
\small
\begin{tabularx}{\textwidth}{C{1cm} X C{2.5cm}}
\toprule
\textbf{\#} & \textbf{Scénario de test} & \textbf{Résultat} \\
\midrule
T1 & Authentification avec identifiants valides/invalides & Validé \\
T2 & Inscription d'un participant à une session déjà complète & Validé \\
T3 & Génération d'un certificat pour un participant déjà certifié (unicité) & Validé \\
T4 & Accès à une route admin avec un jeton formateur (refus attendu) & Validé \\
T5 & Génération de documents PDF (convention, contrat, certificat) & Validé \\
T6 & Signature électronique puis téléchargement du document archivé & Validé \\
T7 & Question posée à l'assistant IA sur un support PDF & Validé \\
\bottomrule
\end{tabularx}
\caption{Extrait du plan de tests fonctionnels}
\end{table}

\subsection{Tests techniques}
\begin{itemize}
  \item Vérification des types : compilation TypeScript sans erreur
    (front-end et back-end) ;
  \item Analyse statique : ESLint sans erreur ;
  \item Build de production : \texttt{nest build} (back-end) et
    \texttt{vite build} (front-end) réussis.
\end{itemize}

% ============================================================
%  CONCLUSION GÉNÉRALE
% ============================================================
\chapter*{Conclusion générale}
\addcontentsline{toc}{chapter}{Conclusion générale}

Ce stage de fin d'études nous a permis de concevoir et de réaliser une
plateforme web complète de gestion des formations pour la STEG. Mené selon
la méthodologie Scrum, le projet a abouti à deux releases couvrant la gestion
des formations (catalogue, sessions, présences, évaluations) puis la
production documentaire automatisée avec signature électronique, le pilotage
par indicateurs et l'exploitation intelligente des supports de formation.

Sur le plan technique, ce projet nous a permis de consolider nos
compétences en développement full-stack (NestJS, React, TypeScript,
MySQL), en modélisation UML, en sécurité applicative (JWT, contrôle d'accès
par rôles) et en intégration d'une solution d'IA appliquée au métier (RAG).

\textbf{Perspectives :}
\begin{itemize}
  \item extension aux formats Word/PowerPoint pour l'assistant IA et ajout
    de l'OCR pour les documents scannés ;
  \item application mobile de pointage des présences ;
  \item gestion des budgets et du remboursement CNFCPP ;
  \item déploiement sur l'infrastructure interne de la STEG avec haute
    disponibilité.
\end{itemize}

% ============================================================
%  BIBLIOGRAPHIE
% ============================================================
\begin{thebibliography}{9}
\addcontentsline{toc}{chapter}{Bibliographie}
\bibitem{scrum} K.~Schwaber, J.~Sutherland, \emph{The Scrum Guide}, 2020.
\bibitem{nestjs} Documentation officielle NestJS, \url{https://docs.nestjs.com}.
\bibitem{react} Documentation officielle React, \url{https://react.dev}.
\bibitem{typeorm} Documentation officielle TypeORM, \url{https://typeorm.io}.
\bibitem{uml} P.~Roques, \emph{UML 2 par la pratique}, Eyrolles, 7\ieme{} édition.
\bibitem{jwt} RFC 7519 : JSON Web Token, IETF.
\bibitem{rag} P.~Lewis et al., \emph{Retrieval-Augmented Generation for
Knowledge-Intensive NLP Tasks}, NeurIPS, 2020.
\end{thebibliography}

% ============================================================
%  ANNEXES
% ============================================================
\appendix
\chapter{Annexes}
\section{Extraits de code}
% Insérez ici 1 ou 2 extraits représentatifs (entité, garde de rôle...)
\lstinputlisting[caption={Exemple : entité Formation (extrait)}]{annexe-formation.ts}

\section{Attestation de stage}
% Déposez le PDF scanné de l'attestation dans le dossier du projet :
\includepdf[pages=-]{attestation-stage.pdf}

\end{document}
